% M1_08: Online Appendix, new short subsection "An informed shareholder" (stake extension).
% INSERT BEFORE: online_appendix.tex line 1611 ("\section{B. Computer-assisted equilibrium certificates}")
% Placement: last subsection of Section A. Renumber "A.12" below to the next free A number at integration.
% Equation displays carry no \tag; number them in the OA sequence at integration.
% NUM placeholders: \NUM{Omega_U} (benchmark value of the full-order bound, in order units),
%   \NUM{Omega_0} (benchmark value of the no-trade bound, in order units).
% New \bibitem requested: none. ref-EdmansGoldsteinJiang2015 is already in the bibliography of the main text;
%   the Online Appendix needs its own \bibitem for that key if it does not have one yet.
% New symbols: \(\omega\) (initial stake, in order units), \(w_j=t_j-t_0\), \(a(q)\), \(\bar e(q)\),
%   \(G_H(q)\), \(G_L(q)\), \(\Omega_U\), \(\Omega_0\), all defined in the text.
%   \(\kappa(\pi)=\min\{m(\pi),1-M(\pi)\}\) follows the notation sheet of the brief.

\subsection{A.12. An informed shareholder}\label{oa-a-stake}

Give the informed investor an initial stake of \(\omega\ge 0\) target
shares, in the same unit as the order \(q\). The stake cost is sunk.
The stake carries no control rights, and the investor still cannot
acquire the target. Its payoff in class \(j\in\{H,L\}\) becomes

\[
\omega V_T+q(V_T-P)-k|q|.
\]

Fix a candidate price and entry schedule. Write \(e(x)\) for entry
at flow \(x\), averaged over preparation costs, and \(\mu_X(x)\) for
the market posterior. Write \(w_j=t_j-t_0\) for the entry gain in
proceeds. Entry raises expected proceeds at each challenger value
above the reserve and leaves proceeds unchanged at a value with no
admissible bid. In the ranked cases \(w_H>w_L\ge 0\), with
\(w_H>0\) when \(\Delta_T>0\). Conditional on
class and flow,

\[
\mathbb E[V_T\mid j,x]=t_0+w_j e(x),\qquad
P(x)=t_0+e(x)[w_L+\Delta_T\mu_X(x)].
\]

The trading residuals keep the form of equation (11) of the main
text. The stake adds a second term. Let
\(D=\{x:\mu_X(x)\ge\tau\}\) be the flow set with high-cost entry,
where \(\tau\) is the threshold belief with gross profit equal to
\(c_H\). Let \(a(q)=\int_D f(x-q)\,dx\) and
\(\bar e(q)=\rho+(1-\rho)a(q)\), the expected entry after order
\(q\) against the fixed schedule. Let

\[
G_H(q)=\int f(x-q)e(x)\Delta_T[1-\mu_X(x)]\,dx,\qquad
G_L(q)=\int f(x-q)e(x)\Delta_T\mu_X(x)\,dx.
\]

Expected payoffs then split into a trading term and a stake term:

\[
\begin{aligned}
U_H(q;\omega)&=\omega[t_0+w_H\bar e(q)]+qG_H(q)-k|q|,\\
U_L(q;\omega)&=\omega[t_0+w_L\bar e(q)]-qG_L(q)-k|q|.
\end{aligned}
\]

For \(q\ne 0\), the derivatives split in the same way:

\[
\begin{aligned}
U_H'(q;\omega)&=
\underbrace{G_H(q)+qG_H'(q)-k\operatorname{sgn}(q)}_{\text{trade and its cost}}
+\underbrace{\omega w_H\bar e'(q)}_{\text{initial stake}},\\
U_L'(q;\omega)&=
\underbrace{-G_L(q)-qG_L'(q)-k\operatorname{sgn}(q)}_{\text{trade and its cost}}
+\underbrace{\omega w_L\bar e'(q)}_{\text{initial stake}}.
\end{aligned}
\]

The constant \(\omega t_0\) drops from each order comparison. The
stake term equals \(\omega\) times the change in expected terminal
proceeds through entry. It does not equal \(\omega\) times a change
in the interim price. A price change alone moves cash between buyer
and seller. An entry change moves \(V_T\) itself, by \(w_j\) per
unit of entry probability.

The entry set stays fixed under a deviation, so the derivative
falls on the density:

\[
a'(q)=\frac1b\left[\int_{D\cap(q,\infty)}f(x-q)\,dx
-\int_{D\cap(-\infty,q)}f(x-q)\,dx\right],
\qquad
|a'(q)|\le\frac{\min\{a(q),1-a(q)\}}b\le\frac1{2b}.
\]

Hence \(|\omega w_j\bar e'(q)|\le\omega(1-\rho)w_j/(2b)\) for each
candidate schedule, with no monotonicity condition on \(D\). The
same step gives \(|G_j'(q)|\le G_j(q)/b\). Against an uninformative
candidate price, \(\mu_X=\pi\) at each flow and entry is constant.
Then \(\bar e'(q)=0\) at each deviation order. The stake adds a
constant and leaves each order comparison at its baseline value.
In particular, the no-trade outcome remains an equilibrium at each
stake. The same point holds against each constant-entry schedule.
Against the full-order schedule, \(D=[x^*,\infty)\) and
\(\bar e'(q)=(1-\rho)f(x^*-q)>0\). A stake then rewards a higher
order in each class with \(w_j>0\). It reinforces the purchase in
class \(H\) and opposes the sale in class \(L\).

Two sufficient bounds keep the baseline outcomes. The residuals
satisfy \(G_H\ge\rho[1-M(\pi)]\Delta_T\) and
\(G_L\ge\rho m(\pi)\Delta_T\). With
\(\kappa(\pi)=\min\{m(\pi),1-M(\pi)\}\), each trading part keeps a
margin of at least \((1-1/b)\rho\kappa(\pi)\Delta_T-k\) in the
correct direction. Since \(w_H>w_L\), the stake bound with \(w_H\)
covers each class. Define

\[
\Omega_U:=
\frac{2b\{(1-1/b)\rho\kappa(\pi)\Delta_T(r_1)-k\}}
{(1-\rho)[t_H(r_1)-t_0(r_1)]}.
\]

For \(0\le\omega<\Omega_U\), the unique trading outcome at \(r_1\)
remains \((q_H,q_L)=(1,-1)\) (analytical). At the benchmark,
\(\Omega_U=\NUM{Omega_U}\) order units. At \(r_0\), each residual is
at most \(\Delta_T(r_0)\), so
\(U_j(q;\omega)-U_j(0;\omega)\le
|q|[\Delta_T(r_0)-k+\omega(1-\rho)w_H(r_0)/(2b)]\). Define

\[
\Omega_0:=
\frac{2b[k-\Delta_T(r_0)]}
{(1-\rho)[t_H(r_0)-t_0(r_0)]}.
\]

For \(0\le\omega<\Omega_0\), no trade at \(r_0\) remains unique
(analytical). At the benchmark, \(\Omega_0=\NUM{Omega_0}\) order
units. A stake below \(\min\{\Omega_0,\Omega_U\}\) therefore keeps
parts (i) and (ii) of Proposition 2 of the main text. Part (iii)
needs no stake bound, because entry there equals \(\rho\) at each
belief and the stake term vanishes. Each bound is sufficient, in
order units. Exact stake cutoffs remain open.

The sign pattern differs from Edmans et al.
(\citeproc{ref-EdmansGoldsteinJiang2015}{2015}, Section II.E). In
that model the real action corrects course in each state. Correct
investment raises the stake value after good news, and correct
disinvestment raises it after bad news. The stake motive is then
symmetric across purchases and sales, and it cancels from the
trading asymmetry. Here entry raises target proceeds in each class,
but the price recruits high-cost entry only after favorable flow. A
sale after bad news lowers the entry probability at the margin and
so lowers the value of the stake. The stake therefore works against
revelation of bad news, and it does not cancel.
